\begin{abstract}
Retail traders who use leveraged contracts for difference (CFDs) on Bitcoin can lose their whole margin in a single day, yet most deep learning work on cryptocurrency is judged only on forecast accuracy. This paper tests whether forecasting risk is more useful than forecasting price. I built a decoupled framework with three parts. The Alpha Agent, an attention-based multivariate Bi-LSTM, forecasts the next day's return. The Beta Agent, two small GRUs, forecasts the size of the next day's move and the chance of a jump. A rule-based Dynamic Risk Engine turns these forecasts into a position, a size, a stop-loss and a leverage level. The system was trained on 2018 to 2023, tuned on 2024, tested on 2025 and then run untouched on nine months of 2026 data that did not exist when it was built. The direction forecasts were no better than a coin flip (50.2\% pooled hit rate), and neither was any simple rule. The Beta Agent's volatility forecast beat naive persistence, but a GARCH(1,1) model put on the same scale was more accurate still. Compared with the same system using fixed risk rules, the Risk Engine cut maximum drawdown from 17.2\% to 4.6\% in 2025 and from 22.8\% to 8.7\% in 2026, both significant under a block bootstrap, without lowering the Sharpe ratio. It did not reduce liquidations, and at 20x leverage it was liquidated twice. The value of the system came from controlling risk, not from predicting direction.
\end{abstract}

\begin{IEEEkeywords}
cryptocurrency, CFD trading, Bi-LSTM, attention, volatility forecasting, risk management, backtesting
\end{IEEEkeywords}
